\section{结果}
\label{sec:results}

\subsection{队列与评估设置}

表~\ref{tab:cohort} 概括了队列及其划分。患者级留出测试集包含来自 4 个中心、
240 名患者的 4,107 张图像。划分在患者层面进行,与训练集无患者重叠;测试集
包含的全部四个中心也出现在训练划分中,因此测试集刻画的是在未见患者(而非
未见机构)上的表现。

测试队列在 F0、F1、F2、F3 各含 60 名患者,而训练与验证遵循队列分布,其中
F2 是患者层面最常见分级。由于 240 名测试患者中每个分期各占 60 名,总体表现
不应被解读为人群筛查场景下的患病率加权估计。图像层面计数并不严格平衡,
因为各患者贡献的图像数量不同。

以下所有结果均使用各配置在预设单一种子~2026 下由验证集选定的检查点。

\subsection{分级表现}

表~\ref{tab:main} 报告了全部五个配置的复合终点及图像级、患者级分项指标。
按 $R_{\mathrm{final}}$ 观察,排序为 E($0.186865$)$<$ A($0.187312$)$<$
D($0.196686$)$<$ B($0.198100$)$<$ C($0.200202$)。因此,SynAP-Fib(E)
取得五个配置中最佳观察 $R_{\mathrm{final}}$。

配置 E 同时取得最低的图像级 COR($0.177908$)与平均绝对误差($0.377857$),
最低的患者最大值 COR($0.200528$)与患者最大值 MAE($0.420558$)。逐指标的
最优分布在各配置间不同:B 取得最高四分类准确率($67.40\%$),C 取得最高
$\pm0.5$ 准确率($70.17\%$);E 取得最低观察 $R_{\mathrm{final}}$。

\subsection{部位与病灶先验的组合}

表~\ref{tab:priors} 将四个拉伸缩放配置排成部位与病灶先验的 $2\times2$
比较。相对于基线,单独引入任一先验均使复合分级指标恶化:$R_{\mathrm{final}}$
从 A 的 $0.187312$ 升至仅含部位先验配置 C 的 $0.200202$,以及仅含病灶先验
配置 D 的 $0.196686$。两个分项 COR 呈同方向变化,因此恶化并不局限于单一
聚合层级。

联合配置 E 未继承这一恶化。它挽回了任一单先验配置下的损失,并取得四种配置
中最佳观察的 $R_{\mathrm{final}}$、图像级 COR 与患者最大值 COR。联合配置
优于仅含部位先验配置的结论得到补充材料中配对 bootstrap 分析的支持
(E~$-$~C 的 95\% CI:$-0.024376$ 至 $-0.002223$)。这是一种非可加的实证
模式:同时提供两个先验的效应,不是两个单先验配置孤立结果所能推断的。

\subsection{分期表现}

图~\ref{fig:stagewise} 在四个临床分期上比较基线与 SynAP-Fib。分期 MAE 由
冻结的逐图像预测事后重算。

分期 MAE 方面,E 在 F0($0.196277$ 对 $0.196508$)、F1($0.411035$ 对
$0.416107$)与 F3($0.401280$ 对 $0.403570$)上低于 A,在 F2($0.516525$
对 $0.514474$)上高于 A。两个配置的最大分期 MAE 均出现在 F2。

分期召回率方面,E 在 F0($81.14\%$ 对 $80.18\%$)与 F1($50.50\%$ 对
$49.50\%$)上高于 A,在 F2 上持平(均为 $63.63\%$),在 F3($72.10\%$ 对
$73.27\%$)上低于 A。两个配置的最低召回率均出现在 F1。

因此,两个指标指认的"最困难分期"并不相同:F2 具有最大的分期 MAE,而 F1
具有最低的召回率。

\subsection{解剖先验输出与计算代价}

除纤维化分数外,联合配置还输出两个附加结果,汇总于表~\ref{tab:auxiliary}。
六类解剖协议下的部位分类在 E 中达到 $66.81\%$ 准确率与 $0.6659$ 的
macro-F1,高于仅含部位先验配置 C 的 $66.40\%$ 与 $0.6623$,交叉熵也更低
($0.9165$ 对 $0.9444$)。与弱边界框的一致程度在 E 中于 4,035 个有效框上
达到 Dice@0.5 $0.3603$ 与 IoU@0.5 $0.2355$,高于仅含病灶先验配置 D 的
$0.3539$ 与 $0.2261$。这些病灶指标量化的是与弱框的一致程度,并非分割真值
表现。

这些附加输出的代价不大。相对于基线,E 增加 1,442,673 个参数($5.9\%$),
乘加运算多 $2.8\%$,单张图像中位延迟由 $3.031$~ms 升至 $3.993$~ms。两条
残差通路的学习门控分布见补充材料。

\subsection{检查点选择诊断}

对配置 A、B、D、E 而言,验证集选定的检查点在该配置的 epoch-120 检查点之上,
测试集表现也更优。配置 C 是唯一例外:其最佳检查点的测试 $R_{\mathrm{final}}$
为 $0.200202$,而 epoch~120 为 $0.191930$。预设选择规则在执行中未做任何
事后替换,报告的排序因而全程使用验证集选定的检查点。

全部预设对比的配对 bootstrap 分析,连同判别、校准、患者中位数与逐中心
结果,见补充材料。
